\scanpage{52}
configuration \deleted{\unsure{icosaèdre|icosaédrale}} \inserted{octaédrale} associée est bien celle de départ\ldots

\uline{Théorème} \quad Soit \(X\) un fibré en droites projectives sur un schéma \(S\), de car. \deleted{\unsure{rési}} résiduelles \(\ne 2\). Considérons le \deleted{\unsure{grp}} schéma en groupes \(G=\underline{\operatorname{Aut}}_S(X)\) (\uline{NB} la donnée de \(X/S\) équivaut à celle de la ``forme'' \(G\) de \(\mathrm{GP}(1)_S\) sur \(S\)), et les ensembles suivants :

a) L'ens \(\mathcal A(X)\) des sous-schémas en groupes \(V\) de \(G\), loc isom. à \((\mathbb D_4)_S\)

b) L'ens \deleted{\unsure{de|a}} \(\mathcal B(X)\) des sous-schémas en groupes \(\deleted{\unsure{\mathfrak S|G}}\Gamma\) de \(G\), loc. isom. à \((\mathfrak S_4)_S\).

c) L'ens des configurations icosaédrales régulières dans \(X\), i.e. des diviseurs étales \(\Delta\) dans \(X/S\), de degré \(6\), munis de un antipodisme \(\underline a_\Delta:\Delta \mathbin{\text{---}} \Delta\) qui \deleted{\unsure{définit un autom.}} \inserted{\unsure{soit|sont} sans pts} fixe sur chaque fibre géom., et tel que \(\underline{\operatorname{Aut}}_{\unsure{S,\rm oct|\rm oct}}^+(\deleted{\unsure{X}}\Delta)\) opère sur \(X\) de façon compatible avec son opération sur \(\Delta\).

On a alors un système transitif de bijections canoniques entre ces trois ensembles, \unsure{obtenues|obtenus|obtenu} ainsi !
% archive sheet number: 51
